%\section{nightline}
\section{Nightline}%FIXME
Die Nightline ist eine telefonische Anlaufstelle von Studierenden für Studierende. Jede kann bei uns anrufen oder eine Mail schreiben, um über alles, was sie gerade beschäftigt, anonym und vertraulich zu reden. Egal ob Ersti oder Doktorandin, ob 18 oder 48 Jahre alt, egal, ob jemand einfach nur kurz etwas loswerden will oder gerade alles über einem zusammenbricht. Die Nightline bietet die Möglichkeit zum offenen Gespräch am späten Abend und nachts, wenn belastende Gefühle und Ängste bestehen, die tagsüber noch gekonnt verdrängt wurden und andere Gesprächspartnerinnen nicht oder nicht mehr erreichbar sind. Wir sind selbst Studierende, befinden uns also in einer ähnlichen Lebenslage wie ihr.

\marginpar{
    \centering{
        \includegraphics[width=4cm]{bilder/nightline_logo.png}
    }
}

Seit Mai 1995 existiert die Nightline in Heidelberg als eingetragener, ehrenamtlich tätiger Verein. Sie wurde von einem Heidelberger Studenten gegründet, der die Idee 1994 nach einem Auslandsaufenthalt aus England „importierte“. Dort gibt es in jeder größeren Unistadt Nightlines, die ihren Service wöchentlich die ganze Nacht lang anbieten. Dies sollte auch in Heidelberg entstehen und nach ausgiebigen Gesprächen mit Expertinnen und einer ersten Schulung begann im Sommersemester '95 der erste Dienst.

Für viele Anrufende ist es sicherlich einfacher, sich erstmal ebenfalls an Studierende zu wenden, da sich diese in einer ähnlichen Lebenssituation befinden. Die Universität und das Studierendenwerk bietet ihren Studierenden zwar eine professionelle Betreuung in der psychotherapeutischen Beratungsstelle, doch ist dazu zunächst eine Terminvereinbarung nötig, weshalb es Studierenden erstmal leichter fällt, bei der Nightline anzurufen. Hier kann man sofort wenn der Schuh drückt, der Kummer überhand nimmt oder man während einer nächtlichen Lerneinheit eine kleine Sinnkrise erleidet anrufen und seinen Problemen Platz machen.

Generell werden wir aus ganz unterschiedlichen Gründen angerufen. Häufige Themen sind zum Beispiel Probleme mit Freunden oder Familie, Stress in der Uni oder Liebeskummer. Außerdem helfen wir auch gerne bei allgemeinen Fragen zur Uni oder zum Unileben weiter.

Einer der Grundpfeiler unserer Arbeit am Telefon ist die Vertraulichkeit. Die Themen, die am Telefon angesprochen werden, bleiben innerhalb der Nightline und kursieren nicht darüber hinaus. Die Anrufende muss ihren Namen nicht nennen und auch die Nightlinerin stellt sich nicht persönlich vor. Dadurch wollen wir einen vorurteilsfreien und neutralen Raum schaffen. Die Anrufende bleibt vollkommen anonym – am Telefon der Nightline ist noch nicht einmal dessen Rufnummer zu sehen. Die bestehende Anonymität kann der Anrufenden helfen, ihr Problem offen auszusprechen ohne das Gefühl zu haben, sich auf irgendeine Art und Weise der Nightlinerin offenbaren zu müssen.

Im Grunde verstehen wir unsere Aufgabe im Zuhören. Uns ist es wichtig vorurteilsfrei, anonym und vertraulich zu sein. Die Briten haben den Leitspruch: „We listen, not lecture!“ Daran halten auch wir uns. Unsere Aufgabe ist es nicht Ratschläge zu erteilen, sondern zuzuhören. Die Nightline wird von zwei Diplompsychologen betreut, die eine Schulung für Nightliner durchführen und mehrmals im Semester Supervisionen abhalten. Sie sind ausgebildete Therapeuten und vermitteln ihre Kenntnisse über Gesprächsführung an uns weiter.

Prinzipiell versuchen wir immer zu helfen und zuzuhören, doch wenn wir merken, dass ein Anrufender mehr braucht, als wir bieten können, oder wir unsere eigenen Grenzen überschritten sehen, leiten wir wenn dies gewünscht wird an professionelle Dienste weiter. Dazu haben wir eine Sammlung von verschiedenen Beratungsstellen und~Nummern angelegt, um beispielsweise an Suchtberatungsstellen, Trauerbegleitungen und fremdsprachige Telefonseelsorgen weiter verweisen zu können.

Die Nightline besteht aus etwa 30 ehrenamtlichen Mitarbeiterinnen aus verschiedensten Fachbereichen. Das Geschlechterverhältnis ist nahezu ausgeglichen, so dass sich in der Nightline das breite Spektrum der Studierendenschaft widerspiegelt. Mitmachen kann jede, die bereit ist, sich ehrenamtlich zu engagieren und einen Teil ihrer Freizeit in den Dienst anderer zu stellen. Ein vorgegebenes Studienfach für Nightliner gibt es nicht.

Die Nightline ist während der Vorlesungszeit täglich zwischen 21 Uhr abends und 2 Uhr nachts unter der Nummer 0\,62\,21 / 18\,47\,08, via Skype unter „nightline.heidelberg“ oder per Mail\footnote{Das Mailsystem findest du unter \url{http://www.nightline-heidelberg.de/}} erreichbar.
